\documentclass[11pt,a4paper]{article}
\usepackage[utf8]{inputenc}
\usepackage[T1]{fontenc}
\usepackage[english]{babel}
\usepackage{amsmath,amssymb,amsthm,mathtools}
\usepackage[margin=2.8cm]{geometry}
\usepackage{enumitem}
\usepackage{hyperref}
\hypersetup{colorlinks=true,linkcolor=black,citecolor=blue,urlcolor=blue}
\emergencystretch=2em
\theoremstyle{plain}
\newtheorem{theorem}{Theorem}[section]
\newtheorem{lemma}[theorem]{Lemma}
\newtheorem{proposition}[theorem]{Proposition}
\newtheorem{corollary}[theorem]{Corollary}
\theoremstyle{definition}
\newtheorem{definition}[theorem]{Definition}
\theoremstyle{remark}
\newtheorem{remark}[theorem]{Remark}
\newcommand{\RP}{\mathbb{RP}}
\newcommand{\PGL}{\mathrm{PGL}}
\newcommand{\rk}{\operatorname{rank}}
\newcommand{\tr}{\operatorname{tr}}
\title{The period spectrum of reflective projective billiards:\\
joins, suspensions, and metric obstructions}
\author{Jorge Lucas Gonz\'alez\thanks{Email: jorge.lucas.glez@gmail.com}\\
\small Independent researcher}
\date{August 2026}
\begin{document}
\maketitle
\begin{abstract}
We determine the complete period spectrum of bounded polyhedral projective
billiards with central facet fields with an open set of primitive periodic trajectories, under the natural
convention that the witness visits all facets. The proof combines projective
joins, rank-one suspensions, and a parity lift of planar centrally projective
billiards. In projective dimension $n$, every period $k\ge n+1$ occurs when $n$
is even, while for odd $n$ exactly the even periods $k\ge n+1$ occur. Metric
examples exist for every pair $k=m+2q$, and positive metric monodromy requires
$k\ge2m$. Thus metric existence is completely determined in odd projective
dimension, and the only remaining branch has odd vector rank, positive monodromy
and even period at least twice the rank.
\end{abstract}

\noindent\textbf{Keywords.} Projective billiards; periodic trajectories; convex
polytopes; joins; reflection factorizations; spherical billiards.

\medskip
\noindent\textbf{2020 Mathematics Subject Classification.}
37C83, 51N15, 52B11, 20F55.

\section{Introduction}
Projective billiards were introduced by Tabachnikov \cite{tab}.  Such a billiard
is \emph{$k$-reflective} when its billiard map has an open set of $k$-periodic
points; Fierobe \cite{fie-ex} gave the first examples, of period three and of
every even period at least four, and Problem~1 of the 2022 problem collection
\cite{opb} asks which further periods occur.  This paper proves the required cyclic constructions and determines the whole
period spectrum.

The basic closure condition for a polyhedral projective billiard is that the
ordered product of its harmonic facet homologies be scalar. Algebraic closure is
not sufficient: the mirrors must bound a proper convex domain and one developed
line must cross the developed facets in strict order. The constructions below
are designed to preserve all three requirements.

The projective join combines negative-monodromy witnesses without losing
boundedness or the one-turn property. A rank-one suspension raises dimension and
period together, and a parity lift embeds a regular polygon as a planar section of
a higher-dimensional polytope. These operations fill the entire period table.
They also expose a sharper metric question, for which we prove a complete
negative-monodromy ray and a sharp lower bound on positive monodromy.

We recall the reflection model and openness criterion, and prove the cyclic atoms and all subsequent constructions here.

For higher-dimensional geometric context, Fierobe's analysis of caustics is
complementary to the polyhedral constructions here \cite{fie-caustics}.

The main contributions of this paper are: the complete period spectrum under
the explicit all-facets witness convention; join, suspension, and parity-lift
constructions that realize the spectrum; and the metric classification that is
proved for the negative-monodromy ray and constrained by the sharp positive
monodromy bound.  The remaining odd-rank, positive-monodromy branch is stated
as open rather than folded into the existence theorem.
\section{Setting}\label{sec:setting}

A \emph{projective billiard} \cite{tab} is a bounded convex polytope
$\Omega\subset\mathbb{R}^n$ together with a \emph{transversal line field}
assigning to each point in the relative interior of a facet of $\partial\Omega$
a line $L(p)$ through $p$ not contained in the facet hyperplane; reflection is
defined by requiring the incoming ray, the outgoing ray, the facet tangent and
$L(p)$ to have cross ratio $-1$ in their impact pencil.

Throughout this paper we restrict to \emph{central facet fields}: on each
facet the transversal lines belong to the pencil through one fixed projective
centre outside the supporting hyperplane. No reduction from arbitrary
transversal fields to this class is asserted. Statements about the period
spectrum use the stated convention on the facets visited by the witness.

\begin{definition}[Fierobe \cite{fie-ex}]\label{def:krefl}
$\Omega$ is \emph{$k$-reflective} if there are a $k$-periodic orbit
$(p_1,\dots,p_k)$ and an \textbf{open} $U_1\times U_2\subset(\partial\Omega)^2$
containing $(p_1,p_2)$ every point of which begins a $k$-periodic orbit.
\end{definition}

The open set is the whole content of the definition, and it is where the difficulty
lives: exhibiting one periodic orbit is easy, exhibiting an open family of initial
impact pairs is not.

\subsection{Reflection as a harmonic homology}

For a covector $L$ and a vector $O$ with $LO\neq0$ put $H=I-2\,OL/(LO)$. Then
$H^2=I$, $HO=-O$, $Hx=x$ for $Lx=0$, and $\det H=-1$.

\begin{proposition}\label{prop:law}
Let $p$ lie on $L$, let $\tau\in\ker L\setminus\{0\}$ be a tangent direction, and
put $\nu=O-p$. The involution induced by $H$
on the pencil at $p$ sends $a\tau+b\nu$ to $a\tau-b\nu$; hence the cross ratio of
(incoming, outgoing; tangent, transversal) is $-1$. So reflection in the side with
transversal field the pencil through $O$ \emph{is} the harmonic homology with axis
the facet hyperplane and centre $O$.
\end{proposition}

\begin{proof}
$H$ fixes $p$ and is an involution, so induces an involution of the pencil at $p$
fixing the two distinct lines $\mathbb{R}\tau$ and $pO$; in the coordinate taking them to $0$
and $\infty$ it is $m\mapsto-m$. Parametrising by $(a:b)$ with direction
$a\tau+b\nu$, the tangent is $(1:0)$, the transversal $(0:1)$, and with
$[ij]=a_ib_j-a_jb_i$ the cross ratio of $(a:b),(a:-b),(1:0),(0:1)$ is
$[13][24]/([14][23])=-1$.
\end{proof}

Consequently, if $H_{j_k}\cdots H_{j_1}$ is scalar then \emph{every} line returns to
itself after the corresponding $k$ reflections. The converse is false: an isolated
trajectory may well close without the product being scalar, and it is precisely the
scalar case that produces an \emph{open set} of closed trajectories, which is what
Definition~\ref{def:krefl} demands. If the product is $cI$, determinants give
$c^{n+1}=(-1)^k$. For odd projective dimension $n$, the left side is positive,
so $k$ must be even. For even $n$, the real scalar is $c=(-1)^k$ and
determinants give no additional obstruction.
For the central fields considered here, an open family also forces scalar
monodromy. The initial lines of a regular family form an open subset of the
Grassmannian of projective lines. Their invariance under the fixed product
is an algebraic incidence condition; vanishing on an open subset forces
it on the entire Grassmannian. A linear map preserving every two-dimensional
vector subspace preserves every one-dimensional subspace, obtained as an
intersection of such subspaces in vector dimension at least three.
Applying this to basis vectors and their pairwise sums shows that the map
is scalar. This necessity concerns an open family; an isolated closed orbit
does not suffice.

\section{An openness criterion}

Everything below rests on this; it assumes no symmetry.

\begin{lemma}[Exit map]\label{lem:exit}
Let $\Omega$ be a convex polytope, $p$ in the relative interior of a facet, $w$ pointing into
$\Omega$. The ray $p+sw$, $s>0$, meets $\partial\Omega$ in exactly one point $E(p,w)$, and
$E$ is continuous wherever $E(p,w)$ lies in the relative interior of a facet.
\end{lemma}

\begin{proof}
$\Omega$ is convex and bounded, so $s^\ast=\sup\{s>0:p+sw\in\Omega\}$ is finite and
$E(p,w)=p+s^\ast w\in\partial\Omega$; uniqueness is convexity. If $E(p,w)$ is
in the relative interior of a facet $F$ then $w$ is not parallel to its
supporting hyperplane, so the intersection determinant is nonzero and Cramer's
rule gives $s^\ast$ near $(p,w)$ as a quotient of polynomials with nonvanishing
denominator. Since the relative interior of $F$ is open in $\partial\Omega$,
nearby data still exit through $F$.
\end{proof}

\begin{remark}
Requiring instead that the intersection determinants with \emph{every} supporting
hyperplane be nonzero fails: $\det=0$ is not an open condition, so such an argument has
a gap. This is why the exit map is isolated.
\end{remark}

\begin{theorem}[Openness criterion]\label{thm:open}
Let $\Omega$ be a convex polytope with facets $F_0,\dots,F_{k-1}$ and transversal
field on $F_i$ the pencil through $O_i$, with $L_iO_i\neq0$. Assume
\textup{(H1)} $H_{k-1}\cdots H_0=cI$, $c\neq0$; \textup{(H2)} there is an orbit
with itinerary $0,1,\dots,k-1$ whose impacts are strictly interior to their facets,
with strictly positive step parameters, each losing facet excluded by a strict
inequality and nonvanishing transversality at each impact; \textup{(H3)} the return
direction is $u_k=\kappa u_0$ with $\kappa>0$. Then $\Omega$ is $k$-reflective; if
the $k$ impacts lie on distinct facets the period is primitive.
\end{theorem}

\begin{proof}
Each step is a line--line intersection followed by a reflection, so the impact
parameters and directions are rational functions of the initial data, with
denominators nonvanishing at the witness by (H2); the vanishing locus being closed,
they are defined and continuous nearby. The finitely many strict inequalities of
(H2) persist, and by Lemma~\ref{lem:exit} each of the $k$ exit maps is continuous
with value in the relative interior of a facet, which is open in $\partial\Omega$; composing,
the itinerary is unchanged for nearby data. By (H1) the product is the identity of
$\PGL_{n+1}$, so \emph{every} line returns to itself: $\ell_k=\ell_0$ as an unoriented
line, and since $\ell_0$ is not the supporting hyperplane of $F_0$ the intersection
$\ell_0\cap F_0$ is a single point, whence $p_k=p_0$. For the orientation, (H1)
only gives $u_k=\kappa u_0$ with $\kappa$ real nonzero --- neither $\pm1$ nor a
priori locally constant --- but $\kappa=(u_k\cdot u_0)/(u_0\cdot u_0)$ is
continuous and nonvanishing and equals a positive number at the witness by (H3), so
$\kappa>0$ on a smaller neighbourhood. Finally, a proper period $d\mid k$, $0<d<k$,
would force $p_d=p_0$; but $p_d$ lies on facet $F_d$ and $p_0$ on $F_0$, which are
distinct facets, and the impacts are interior to them, so $p_d\neq p_0$.
\end{proof}

Hypotheses (H2)--(H3) are finitely many strict inequalities and (H1) one identity,
so they can be certified exactly for explicit data.  We use this repeatedly
below.
\section{The minimal negative-monodromy atom}

\begin{theorem}[The opposite-vertex simplex]\label{thm:opposite-simplex}
For every $n\ge2$, every projective $n$-simplex admits a projective billiard
which is $(n+1)$-reflective with primitive period $n+1$.  On each facet, take as
transversal field the pencil of lines through the opposite vertex.  In fact, for
each ordering of the $n+1$ facets there is a strict periodic witness with that
itinerary.
\end{theorem}

\begin{proof}
By projective invariance it is enough to use the standard simplex.  Let
$V=\mathbb R^{n+1}$ with basis $e_0,\ldots,e_n$, let
\[
 C=\{x=(x_0,\ldots,x_n):x_i>0\text{ for every }i\},
 \qquad \Delta=\mathbb P(C),
\]
and use the bounded affine section $\sum_i x_i=1$.  The facet $A_i$ has equation
$x_i=0$ and its opposite vertex is $E_i=[e_i]$.  The harmonic homology with axis
$A_i$ and centre $E_i$ is
\[
 H_i=I-2e_i e_i^*,
\]
so it changes the sign of the $i$-th coordinate and fixes all the others.  Hence
the homologies commute and
\[
             H_nH_{n-1}\cdots H_0=-I_V.                 \tag{1}
\]

It remains only to certify a strict itinerary.  Choose real numbers
\[
        \lambda_0>\lambda_1>\cdots>\lambda_n>0
\]
and define a linear injection $f:\mathbb R^2\to V$ by
\[
        f(s,t)=(s-\lambda_0t,\ldots,s-\lambda_nt).
\]
Along the oriented half-circle $(s,t)=(\cos\theta,\sin\theta)$,
$0\le\theta\le\pi$, coordinate $i$ vanishes once, at
\[
        \theta_i=\arctan(1/\lambda_i),
\]
and
$0<\theta_0<\theta_1<\cdots<\theta_n<\pi/2$.  Thus the developed projective
line passes successively through the chambers
\[
 C,\quad H_0C,\quad H_1H_0C,\quad\ldots,\quad
 H_n\cdots H_0C=-C.
\]
Between two successive crossings every coordinate has the strict sign prescribed
by the current chamber.  At the $i$-th crossing exactly one coordinate vanishes,
so the point lies in the relative interior of the developed copy of $A_i$ and the
crossing is transverse.  Consequently the folded ray has $A_i$ as its unique first
hit at every step; all admissibility conditions are strict.

The half-circle ends at $(-1,0)$, and $f(-1,0)=-f(1,0)$.  Combining this with
(1), folding returns the endpoint and the oriented tangent to their initial values
with positive multiplier.  The openness criterion therefore supplies an open set
of $(n+1)$-periodic trajectories.  Since all $n+1$ facets are distinct, the period
is primitive.  Permuting the coordinates gives any desired facet ordering.
\end{proof}

\begin{remark}
For $n=2$ this is precisely the right-spherical triangle: every transversal on a
side passes through the opposite vertex.  The proof shows that the natural object
in arbitrary dimension is the projectivisation of the all-right spherical orthant.
It attains the general rank lower bound $k\ge n+1$ exactly.
\end{remark}

\section{Joins, suspensions and filled regions}

\begin{theorem}[Join construction]\label{thm:join}
For $r=1,2$, let $V_r$ be a real vector space, let $C_r\subset V_r$ be an open
pointed polyhedral cone whose projectivisation $\Omega_r=\mathbb P(C_r)$ is a
bounded convex polytope, and let
\[
 H^{(r)}_0,\ldots,H^{(r)}_{k_r-1}\in\mathrm{GL}(V_r)
\]
be the harmonic homologies of a strict billiard orbit, in itinerary order. Assume
that their product is $-I_{V_r}$ and that the developed orbit goes from a lift $p_r$
to $-p_r$ while traversing its projective line exactly once.  Then the projective
join
\[
 \Omega_1*\Omega_2:=\mathbb P(C_1\oplus C_2)
       \subset\mathbb P(V_1\oplus V_2)
\]
with the block homologies
\[
 H^{(1)}_a\oplus I_{V_2}\quad(0\le a<k_1),\qquad
 I_{V_1}\oplus H^{(2)}_b\quad(0\le b<k_2)
\]
is $(k_1+k_2)$-reflective.  Its dimension is
$\dim\Omega_1+\dim\Omega_2+1$ and it has $k_1+k_2$ facets.  If each factor orbit
visits distinct facets, the resulting period is primitive.  Moreover, the
constructed witness also traverses its developed projective line exactly once, so
the construction may be iterated.
\end{theorem}

\begin{proof}
Boundedness is immediate at the level of cones.  Choose linear forms $\ell_r$ that
are positive on $\overline C_r\setminus\{0\}$.  Then
$\ell_1\oplus\ell_2$ is positive on
$\overline{C_1\oplus C_2}\setminus\{0\}$, and its level-one section is compact.
The facets of $C_1\oplus C_2$ are precisely
$F^{(1)}_a\oplus C_2$ and $C_1\oplus F^{(2)}_b$.  Moreover, if
$H=I-2O L/(LO)$ is a factor homology, then
\[
 H\oplus I=I-2(O,0)(L,0)/(LO),
\]
so the displayed block maps give exactly the harmonic reflection laws on these
facets.  Their product, in the concatenated order, is
$(-I_{V_1})\oplus(-I_{V_2})=-I_{V_1\oplus V_2}$.

It remains to produce the strict itinerary; this is the only non-formal point.
Unfold the two given witnesses.  Thus there are two-dimensional spaces $E_r$,
linear injections $f_r:E_r\to V_r$, and oriented half-circles from $e_r$ to $-e_r$
in $E_r\setminus\{0\}$ whose projectivisations cross, in order, the developed
facets of the successive cones
\[
 C_r,\quad G^{(r)}_1C_r,\quad\ldots,\quad
 G^{(r)}_{k_r}C_r=-C_r,
\]
where $G^{(r)}_a$ is the appropriate prefix product.  Strictness says that between
two consecutive crossings the developed line lies in the interior of the
corresponding cone, and at a crossing it lies in the relative interior of just the
prescribed facet.

Use the first developed projective line as parameter circle.  After its last
crossing choose a closed interval $J$ strictly inside its final chamber interval
$-C_1$.  On the second parameter circle choose the complementary cut inside its
initial/final chamber, so that the compact arc $K$ contains all its facet crossings.
The action of $\mathrm{PGL}_2(\mathbb R)$ on the projective line is triply
transitive: there is an orientation-preserving projectivity
$\psi:\mathbb P(E_1)\to\mathbb P(E_2)$ which sends the common starting parameter to
the second starting parameter and sends the two endpoints of $J$ to the two
endpoints of $K$.  Consequently every second-factor crossing has its preimage in
$J$, in the original order.  Before $J$ the second component remains strictly in
$C_2$; throughout $J$ the first component remains strictly in $-C_1$.

Let $A:E_1\to E_2$ induce $\psi$, choosing its overall sign so that the initial
lifts lie in $C_1$ and $C_2$, and put
\[
       F(e)=f_1(e)\oplus f_2(Ae).
\]
The projective line $\mathbb P(F(E_1))$ first crosses all developed facets of the
first factor and then all those of the second.  Between crossings both components
lie strictly in their current developed cones, hence their direct sum lies strictly
in the corresponding developed join.  At a first-factor crossing the first
component lies in the relative interior of its prescribed facet and the second is
interior, and conversely for a second-factor crossing.  Thus every prescribed facet
is the unique first hit: this gives all the strict inequalities and transversalities
in \textup{(H2)}.

Finally the common half-circle ends at $-e$.  Since $F(-e)=-F(e)$ and the total
prefix is $-I$, folding sends the endpoint back to the starting point and sends the
terminal oriented tangent to the initial oriented tangent with positive multiplier.
This is \textup{(H3)}.  The openness criterion now gives an open set of
$(k_1+k_2)$-periodic orbits.  If the factor facets are all distinct, so are the join
facets, and the period is primitive.  Since the parameter runs along the same
oriented half-circle from $e$ to $-e$ without a second turn, the output witness has
the one-turn property claimed in the statement.  More explicitly, the quotient
of that half-circle obtained by identifying its endpoints is mapped
homeomorphically onto $\mathbb P(E_1)$, and the linear injection $F$ induces a
projective isomorphism $\mathbb P(E_1)\to\mathbb P(F(E_1))$.  Hence the developed
line has degree one, not merely odd degree; no metric length estimate or chord
budget is involved.
\end{proof}

\begin{proposition}[Rank-one suspension]
\label{prop:rank-one-suspension}
Let a strict one-turn witness in $\mathbb P(V)$ have period $k$ and product
$-I_V$.  Then the pyramid over its domain in $\mathbb P(V\oplus\mathbb R)$ has a
strict one-turn witness of primitive period $k+1$ and product
$-I_{V\oplus\mathbb R}$.  Thus
\[
                         (n,k)\longmapsto(n+1,k+1).
\]
The old lifted mirror group embeds in the new one.
\end{proposition}

\begin{proof}
Write the old developed lift as $f(e)$ along an oriented half-circle from $e_0$
to $-e_0$.  Choose a nonzero linear form $h$ on the two-dimensional parameter
plane whose unique zero on that half-circle occurs strictly after the last old
facet crossing and before the endpoint.  Choose its sign so that $h(e_0)>0$.
Then
\[
                         F(e)=f(e)\oplus h(e)
\]
starts in $C\oplus\mathbb R_{>0}$.  During every old crossing the last component
is positive.  After the last one, the old component lies strictly in $-C$; when
$h$ vanishes, $F(e)$ crosses the new base facet in its relative interior and
then lies in $-C\oplus\mathbb R_{<0}=-(C\oplus\mathbb R_{>0})$.

The old reflections lift as $H_i\oplus1$, and the new base reflection is
$J=I_V\oplus(-1)$.  Hence the total product in itinerary order is
\[
        J\bigl((-I_V)\oplus1\bigr)=-I_{V\oplus\mathbb R}.
\]
All crossings are strict, $F$ is injective, and the same parameter half-circle
has degree one on $\mathbb P(F(E))$, so closure, positive tangent return, and the
one-turn property follow exactly as in Theorem~\ref{thm:join}.  The new base is
distinct from all old facets, proving primitivity.  Finally, the maps
$H_i\oplus1$ contain an isomorphic copy of the old lifted group.
\end{proof}



\begin{proposition}[Join construction for positive monodromy]
\label{prop:join-positive}
Assume instead that two strict factor witnesses have product $+I$ and that each
developed lift makes exactly one full turn on its two-dimensional vector plane
before returning to its initial vector.  Then the same block construction gives a
strict $(k_1+k_2)$-reflective join with product $+I$.  The output witness again
makes exactly one full vector turn, so this construction may also be iterated.
\end{proposition}

\begin{proof}
Use oriented parameter circles $e_r(\theta)$, $0\le\theta\le2\pi$, for the two
factor lifts $f_r(e_r(\theta))$.  Cut each circle at a parameter lying strictly
inside a chamber interval and whose antipode is not a crossing; this is possible
because the crossing sets are finite.  Choose an orientation-preserving linear
isomorphism $A:E_1\to E_2$ which sends the first cut ray to the second with
positive sign.  Among such maps, choose $A$ generically so that no crossing time
of the first factor coincides with a crossing time of the reparametrized second
factor.  This is possible because there are only finitely many forbidden
equalities, each a proper condition on $A$.

Put $F(e)=f_1(e)\oplus f_2(Ae)$ and order the union of the two crossing sets around
the parameter circle.  Between two consecutive crossings each component lies in
the interior of its current developed cone.  At a crossing exactly one component
lies in the relative interior of its prescribed facet and the other remains
interior.  Thus the direct sum has precisely the merged strict itinerary.  The
block product is $I\oplus I=I$.  Since $e(2\pi)=e(0)$, both the endpoint and the
oriented tangent return with positive multiplier.  Finally, $F$ is injective and
maps one vector-circle turn to one vector-circle turn in $F(E_1)$, proving the
iterability assertion.
\end{proof}

\begin{lemma}[Every even central polygon has one full turn]
\label{lem:even-cyclic-full-turn}
Let $p\ge4$ be even. The centrally projective billiard on the regular $p$-gon
has product $+I_3$, and its midpoint witness makes exactly one full turn on its
developed vector plane.
\end{lemma}

\begin{proof}
Write $\alpha=\pi/p$, $v=\cos\alpha$, and
$\phi_i=(2i+1)\alpha$. At the common-centre choice the reflections have the
explicit form
\begin{equation}\label{eq:even-central-H}
 H_i(x,y,z)=\left(x,y,-z+\frac2v
       (x\cos\phi_i+y\sin\phi_i)\right).
\end{equation}
The midpoint impacts are $M_i=(v\cos\phi_i,v\sin\phi_i,1)$. Let
$G_i=H_{i-1}\cdots H_0$ and $X_i=G_i^{-1}M_i$. The first developed line is
spanned by $X_0=M_0$ and $X_1=H_0M_1$, and direct substitution in
\eqref{eq:even-central-H} puts both in
\[
                         E=\{z=x/v^2\}.
\]
Every prefix preserves $x,y$, and substitution gives
\begin{equation}\label{eq:even-central-f}
 X_i=(v\cos\phi_i,v\sin\phi_i,\cos\phi_i/v)
     =f(\cos\phi_i,\sin\phi_i),\qquad
 f(a,b)=(va,vb,a/v).
\end{equation}
Thus the oriented developed points run from angle $\alpha$ to
$2\pi+\alpha$ in $p$ steps of size $2\alpha\in(0,\pi)$: exactly one full
vector turn in $E$.

For the product, put $a_i=(2/v)(\cos\phi_i,\sin\phi_i)$. Each $H_i$ has
block form $\left(\begin{smallmatrix}I_2&0\\a_i^T&-1\end{smallmatrix}\right)$.
Their product has diagonal blocks $I_2$ and $(-1)^p=1$; the remaining row is,
up to overall sign, the alternating sum of the $a_i$. This sum vanishes because
\[
 \sum_{i=0}^{p-1}(-1)^i e^{\mathrm i\phi_i}
 =e^{\mathrm i\alpha}\frac{1-(-e^{2\mathrm i\alpha})^p}
                         {1+e^{2\mathrm i\alpha}}=0.
\]
Hence the product is $+I_3$.
\end{proof}

\begin{lemma}[A rational positive atom of rank four]
\label{lem:rational-positive-cube}
There is a strict primitive six-periodic projective billiard in a cube in
$\mathbb {RP}^3$ whose canonical lifted product is $+I_4$, whose developed lift
makes exactly one full vector turn, and whose mirror group is infinite.
\end{lemma}

\begin{proof}
Use homogeneous coordinates $(x,y,z,t)$ and the cube
\[
 |x|<z,\qquad |y|<z,\qquad |t|<z.
\]
With the convention that the interior is $L_i<0$, take
\[
\begin{array}{lll}
L_0=(1,0,-1,0),&L_1=(0,1,-1,0),&L_2=(-1,0,-1,0),\\
L_3=(0,-1,-1,0),&L_4=(0,0,-1,1),&L_5=(0,0,-1,-1),
\end{array}
\]
and centres
\[
 \begin{gathered}
 O_i=(0,0,1,c_i)\quad(0\le i<4),\\
 (c_0,c_1,c_2,c_3)=\frac15(-6,-1,4,9),\qquad
 O_4=O_5=(0,0,0,1).
 \end{gathered}
\]
Set $H_i=I-2O_iL_i/(L_iO_i)$.  Direct rational multiplication in the
itinerary order $(1,2,3,4,5,0)$ gives
\[
                         H_0H_5H_4H_3H_2H_1=I_4.
\]
In the affine chart $z=1$, launch from
\[
 X=(2/5,3/10,1,-1/4),\qquad U=(-1/4,1/3,0,1/10).
\]
Exact ray tracing visits
\[
                         1,2,3,4,5,0,1
\]
with a unique positive first-hit time at every step and with every other facet
inequality strict.  After six reflections both the point and the tangent return
exactly, with tangent multiplier one.  Thus the openness criterion applies, and
the six distinct facets make the period primitive.

For completeness, in the basis $(X,U)$ of the developed plane the seven impact
vectors are
\[
 \begin{gathered}
 (1,21/10),\ (-5/2,0),\ (-2,-6/5),\ (-160/91,-118/91),\\
 (40/91,-198/91),\ (1,-12/5),\ (1,21/10).
 \end{gathered}
\]
Every consecutive determinant is positive.  Their quadrant pattern is
$I$, negative axis, $III$, $III$, $IV$, $IV$, $I$; hence the developed lift
makes exactly one, rather than several, positive vector turns.  Finally,
$H_5H_4$ is a nonidentity unipotent with $(H_5H_4-I)^2=0$, so it has infinite
order.  All these equalities and all strict inequalities are independently
checked over $\mathbb Q$ by
\texttt{verificadores/}\allowbreak\texttt{prism\_positive\_rp3.py}, which
exits with code zero.
\end{proof}

\begin{lemma}[A rational positive atom of rank five]
\label{lem:rational-positive-hypercube}
There is a strict primitive eight-periodic projective billiard in a hypercube in
$\mathbb {RP}^4$ with canonical lifted product $+I_5$, exactly one full vector
turn, and infinite mirror group.
\end{lemma}

\begin{proof}
Add a coordinate $s$ to the cube data of
Lemma~\ref{lem:rational-positive-cube}.  For its six old facets and centres put
\[
 \begin{gathered}
 \widehat L_i=(L_i,0),\qquad \widehat O_i=(O_i,d_i),\\
 (d_0,\ldots,d_5)=\left(-3,-\frac{17}{10},-\frac25,
                         \frac9{10},\frac3{10},\frac3{10}\right).
 \end{gathered}
\]
Add the facets and centres
\[
 \widehat L_6=(0,0,-1,0,1),\quad
 \widehat L_7=(0,0,-1,0,-1),\qquad
 \widehat O_6=\widehat O_7=(0,0,0,0,1).
\]
They bound the hypercube $|x|,|y|,|t|,|s|<z$.  With
$\widehat H_i=I-2\widehat O_i\widehat L_i/(\widehat L_i\widehat O_i)$,
exact multiplication in itinerary order $(1,2,3,4,5,6,7,0)$ gives $+I_5$.
The rational launch
\[
 X=(4/11,8/11,1,2/3,-8/15),\qquad
 U=(-4/5,1/3,0,-2/15,1/2)
\]
has strict first-hit itinerary
\[
                         1,2,3,4,5,6,7,0,1
\]
and returns exactly with tangent multiplier one.  The eight successive
determinants of its developed impact vectors are
\[
 \frac{117}{100},\ \frac{1771}{1180},\ \frac{50765}{441438},\
 \frac{825}{1247},\ \frac{680845}{15679778},\ \frac{3300}{6287},\
 \frac{455}{25148},\ \frac{71}{44};
\]
all are positive.  Their quadrant pattern
$I,I,III,III,IV,IV,IV,IV,I$ certifies exactly one vector turn.  The two new
cap reflections again have a nontrivial unipotent product.  The independent
exact checker \texttt{verificadores/}\allowbreak
\texttt{hypercube\_positive\_rp4.py} verifies every equality and strict
inequality over $\mathbb Q$ and exits with code zero.
\end{proof}

\begin{theorem}[The central hypercube family]
\label{thm:central-hypercube}
For every even $n\ge2$, the projective $n$-cube admits a strict primitive open
family of period $2n$ with canonical lifted product $+I_{n+1}$.  Its witness
makes exactly one full vector turn and its mirror group is infinite.
\end{theorem}

\begin{proof}
In homogeneous coordinates $(x_0,\ldots,x_{n-1},z)$ take
\[
 C=\{\,|x_j|<z:0\le j<n\,\},\qquad O=(0,\ldots,0,1),
\]
and give all facets the common centre $O$.  Their covectors are
$L_{j,\sigma}=\sigma x_j^*-z^*$, $\sigma=\pm1$.  Order the facets cyclically as
\[
 L_{0,+},L_{1,+},\ldots,L_{n-1,+},
 L_{0,-},L_{1,-},\ldots,L_{n-1,-}.
\]
Put $H_i=I-2OL_i/(L_iO)$.  Relative to
$V=\ker z^*\oplus\mathbb RO$, every $H_i$ has block form
\[
                    H_i=\begin{pmatrix}I_n&0\\ b_i&-1\end{pmatrix}.
\]
For an even number of factors their product is identity precisely when the
alternating sum of the rows $b_i$ vanishes.  Opposite facets have opposite rows
and their positions differ by $n$.  Since $n$ is even, they have the same parity
and cancel pairwise.  Hence the displayed word has product $+I_{n+1}$.

It remains to realise the word dynamically.  Let $E$ be the affine two-plane
through $O$ parametrised by
\[
 x_j=a\cos\frac{j\pi}{n}+b\sin\frac{j\pi}{n}\qquad(0\le j<n).
\]
Its section by the cube is
\[
 \left|a\cos\frac{j\pi}{n}+b\sin\frac{j\pi}{n}\right|<1
 \qquad(0\le j<n),
\]
the regular $2n$-gon circumscribed about the unit circle.  Every $H_i$ preserves
the homogeneous span of $E$, because it fixes its facet axis and its centre lies
in $E$; its restriction is exactly the centrally projective reflection in the
corresponding side of that regular polygon.  The midpoint witness of
Lemma~\ref{lem:even-cyclic-full-turn} therefore visits the $2n$ cube facets in
the displayed cyclic order.  Every impact is in the relative interior of its
cube facet, since all other section inequalities are strict there.  The same
lemma gives one full vector turn and positive return multiplier, so the openness
criterion applies in the ambient dimension and the distinct facets make the
period primitive.

Finally, for two distinct facets the product of their common-centre homologies is
a nonidentity unipotent of the form
$\left(\begin{smallmatrix}I_n&0\\ b_j-b_i&1\end{smallmatrix}\right)$, hence has
infinite order.
\end{proof}

\begin{corollary}[A positive-monodromy ray in ranks divisible by three]
\label{cor:join-positive-ray}
Let $r\ge1$ and $n=3r-1$.  For every even $k\ge4r$ there is a bounded convex
projective billiard in $\mathbb R^n$ with a primitive open family of period $k$
and canonical lifted product $+I_{n+1}$.  Its mirror group is infinite.
\end{corollary}

\begin{proof}
Join $r$ even central planar atoms.  Take $r-1$ of period four and one of period
$4+(k-4r)$, which is again even and at least four.  Their vector ranks add to
$3r=n+1$, their periods add to $k$, and Proposition~\ref{prop:join-positive}
supplies strictness, primitivity, closure, and one full turn.

For infinitude, at $\lambda=0$ formula~\eqref{eq:even-central-H} writes each
planar generator as
\[
 H_i=\begin{pmatrix}I_2&0\\ a_i^T&-1\end{pmatrix},
 \qquad a_i=\frac2v(\cos\phi_i,\sin\phi_i).
\]
For $i\ne j$,
\[
 H_jH_i=\begin{pmatrix}I_2&0\\ (a_j-a_i)^T&1\end{pmatrix}
\]
is a nontrivial unipotent and therefore has infinite order.  Each joined group
contains a factor group by Proposition~\ref{prop:join-infinite}, so it is infinite.
\end{proof}

\begin{proposition}[Infinite mirror groups survive joins]\label{prop:join-infinite}
Use the involutive linear lifts of all factor homologies.  The lifted reflection
group of either factor embeds in the projective mirror group of the join by
\[
              g\longmapsto [g\oplus I].
\]
Consequently, if either factor has infinite mirror group, then so does the join.
\end{proposition}

\begin{proof}
Every generator $H^{(1)}_a\oplus I$ belongs to the joined group, so the displayed
map is a homomorphism on the linear group generated by the chosen involutive lifts.
If $g\oplus I$ is projectively scalar, say $g\oplus I=cI$, then its second block
gives $c=1$ and its first block gives $g=I$.  Thus the map has trivial kernel.
An infinite projective factor group has an infinite lifted group, proving the last
claim.
\end{proof}

\begin{corollary}[Positive monodromy in the semigroup of ranks $3$ and $4$]
\label{cor:positive-rank-semigroup}
Put
\[
                    K(m)=2\left\lceil\frac{2m}{3}\right\rceil .
\]
For every $m\ge6$, rank $m$ admits an infinite-mirror-group positive-monodromy
witness for every even period $k\ge K(m)$.  In addition, rank $3$ admits every
even period at least $4$, rank $4$ admits period $6$, and rank $5$ admits period
$8$.
\end{corollary}

\begin{proof}
The positive atoms have rank and period $(3,p)$ for every even $p\ge4$ by
Lemma~\ref{lem:even-cyclic-full-turn}, $(4,6)$ by
Lemma~\ref{lem:rational-positive-cube}, and $(5,8)$ by
Lemma~\ref{lem:rational-positive-hypercube}.  Apply
Proposition~\ref{prop:join-positive}.  For $m\ge6$, according as its residue
modulo three is $0,1,2$, use respectively only rank-three atoms, one rank-four
atom, or one rank-five atom, filling the remaining rank with rank-three atoms.
The base period is the least even integer not smaller than $4m/3$, namely
$K(m)$, and every representation contains a planar atom.  Increasing that
atom's period by an arbitrary even amount gives every even $k\ge K(m)$.  Each
join contains an infinite factor group, so its mirror group remains infinite.
\end{proof}



\begin{lemma}[Every odd cyclic atom has one turn]
\label{lem:odd-cyclic-one-turn}
Let $p\ge3$ be odd and put $\alpha=\pi/p$, $u=\sin\alpha$,
$v=\cos\alpha$, and $\phi_i=(2i+1)\alpha$. Define
\[
 \lambda=v(2v+1),\qquad Q=\operatorname{diag}(1,1,\lambda v),
\]
\[
 O_i=(-\lambda\cos\phi_i,-\lambda\sin\phi_i,1)^T,\qquad
 L_i=(-\cos\phi_i,-\sin\phi_i,v),\qquad
 H_i=I-\frac{2O_iL_i}{\lambda+v}.
\]
Let $M_i=(v\cos\phi_i,v\sin\phi_i,1)^T$ be the midpoint impacts. The product
of the $H_i$ is $-I_3$, and the midpoint witness unfolds from a lift $x$ to $-x$
in exactly one turn.
\end{lemma}

\begin{proof}
The form $Q$ is positive definite and $QO_i=\lambda L_i^T$, so each $H_i$ is a
$Q$-reflection. Let $R$ be rotation by $2\alpha$ in the first two coordinates.
Then $H_i=R^iH_0R^{-i}$ and $M_i=R^iM_0$. With $W=H_0R^{-1}$, direct
multiplication gives
\[
 \operatorname{tr}W=\frac{\lambda+2v\cos(2\alpha)-v}{\lambda+v}=2v-1,
 \qquad \det W=-1.
\]
An orthogonal map of a positive three-dimensional space with determinant $-1$
has an eigenvalue $-1$. If its other two eigenvalues were real, the trace would be
$1$ or $-3$; here $0\le 2v-1<1$, so they are a conjugate pair. The trace then
gives $\cos\theta=v=\cos\alpha$, hence the eigenvalues are
$-1,e^{i\alpha},e^{-i\alpha}$. The telescoping product identity
\[
 H_{p-1}\cdots H_0=R^{-1}W^pR
\]
therefore gives $H_{p-1}\cdots H_0=-I_3$. In particular, for
$G_i=H_{i-1}\cdots H_0$ and $X_i=G_i^{-1}M_i$, one has $X_0=M_0$ and
$X_p=-M_0$.

Cyclic symmetry reduces the $Q$-angle of every developed segment to that between
$M_0$ and $H_0M_1$. The point $M_0$ lies on the mirror $L_0$, so self-adjointness
of $H_0$ gives
\[
 \langle M_0,H_0M_1\rangle_Q
 =\langle H_0M_0,M_1\rangle_Q
 =\langle M_0,M_1\rangle_Q.
\]
Using $M_0=(v^2,uv,1)$ and $M_1=(v\cos3\alpha,v\sin3\alpha,1)$ yields
\[
 \|M_0\|_Q^2=2v^2(v+1),\qquad
 \langle M_0,M_1\rangle_Q=2v^3(v+1).
\]
Thus each developed arc has cosine $v=\cos\alpha$ and length $\alpha$. The
original bounded polygon lies in an open affine hemisphere; unfolding by the
$Q$-isometries keeps each segment in an open hemisphere, so these are the minor
arcs. Their total length is $p\alpha=\pi$, the spherical distance between the
antipodal endpoints $X_0$ and $X_p$. The path is therefore a minimizing
semicircle and makes exactly one vector turn.
\end{proof}



\begin{corollary}[The full negative-monodromy ray]
\label{cor:join-full-negative-ray}
For every $n\ge2$ and every $q\ge0$ there is a bounded convex
projective billiard in $\mathbb R^n$ with a primitive open family of period
\begin{equation}\label{eq:full-negative-ray}
                         k=n+1+2q.
\end{equation}
Thus, in every odd dimension $n\ge3$, every period allowed by the determinant
and rank obstructions is attained.  In every even dimension $n\ge2$, every odd
period at least $n+1$ is attained.
\end{corollary}

\begin{proof}
Start with the planar atom of Lemma~\ref{lem:odd-cyclic-one-turn} having period
$p=3+2q$.  Apply the rank-one suspension of
Proposition~\ref{prop:rank-one-suspension} exactly $n-2$ times.  The resulting
dimension is $n$ and the period is
\[
                    (3+2q)+(n-2)=n+1+2q.
\]
Every suspension preserves product $-I$, one turn, strictness, and primitivity.
Finally, the general rank bound is $k\ge n+1$, while in odd dimension the
determinant obstruction forces $k$ to be even; these are exactly the values in
\eqref{eq:full-negative-ray}.
\end{proof}



\begin{remark}[Three independent exact certificates]
The cases $(5,10)$, $(6,11)$ and $(9,14)$ have also been checked without appealing
to the topological synchronisation argument.  The scripts
\texttt{verificadores/join\_rp5.py}, \texttt{join\_rp6.py} and
\texttt{join\_rp9.py} respectively certify over $\mathbb Q$ the scalar product,
boundedness, all strict first-impact inequalities, relative-interior impacts,
primitive period, exact closure and positive return multiplier $\kappa=1$.

For example, in the $(9,14)$ case, on the parameter line of each copy of the
rational $B_5$ witness, with coordinates $e=(a,b)$ and lift
$f(a,b)=aX_0+bU_0$, take
\[
 A=\begin{pmatrix}-200&40000\\-1&100\end{pmatrix},
 \qquad e_*=(1,1/100).
\]
Then $\det A=20000>0$ and $Ae_*=(200,0)$, so the joined line
$F(e)=f(e)\oplus f(Ae)$ starts in the relative interior of the second copy of
facet~$4$.  Exact rational ray tracing gives the fourteen first impacts
\[
5,6,0,1,2,3,4,\quad 5',6',0',1',2',3',4',
\]
all strict, and returns with $p_{14}=p_0$ and tangent multiplier $\kappa=1$.
Every equality and inequality is certified over $\mathbb Q$ by
\texttt{verificadores/join\_rp9.py}, which exits with code zero only if all checks
pass.  Thus the existence of the $(9,14)$ example is independently secured even
before using the join theorem inductively.
\end{remark}

\begin{remark}[Why the one-turn hypothesis is present]
The relation $\prod H_i=-I$ by itself gives only
$L\equiv\pi\pmod{2\pi}$ for a developed spherical orbit; it does not distinguish
$\pi$ from $3\pi,5\pi,\ldots$.  The synchronisation above uses a single traversal
of the parameter circle.  Thus Theorem~\ref{thm:join} does not claim that arbitrary
negative-monodromy witnesses can be joined.  The short-branch hypothesis is an
independent certified input in the corollary.
\end{remark}

\section{A parity lift and the complete existence table}
\begin{proposition}[Central criterion]
\label{prop:common-centre-product}
Let a full-dimensional bounded affine $n$-polytope containing the origin have
genuine facets
\[
                         N_i(w)<1\qquad(0\le i<k),
\]
and give every facet the transversal pencil through the origin.  If $k$ is even,
the canonical lifted product in facet order is $+I_{n+1}$ if and only if
\[
                         \sum_{i=0}^{k-1}(-1)^iN_i=0.       \tag{1}
\]
If $k$ is odd the product is not scalar.  In particular, a bounded common-centre
example with scalar product has $k\ge n+2$, and this bound is sharp whenever $k$
is even.
\end{proposition}

\begin{proof}
Homogenise with a coordinate $z$, put $O=(0,1)\in\mathbb R^n\oplus\mathbb Rz$,
and write $L_i=N_i-z^*$.  Since $L_iO=-1$, the harmonic homology is
\[
 H_i=I-\frac{2OL_i}{L_iO}
    =\begin{pmatrix}I_n&0\\2N_i&-1\end{pmatrix}.           \tag{2}
\]
Multiplying (2) shows that an even product has diagonal blocks $I_n,1$ and
lower-left block equal, up to an irrelevant overall sign, to twice the alternating
sum in (1).  An odd product has diagonal blocks $I_n,-1$ and hence cannot be
scalar.  This proves the criterion.

For the bound, boundedness implies that the $N_i$ span $(\mathbb R^n)^*$ and that
they have a linear dependence with all coefficients positive.  Relation (1) is
independent of that positive dependence, since its coefficients have both signs.
Thus the nullity of the $n\times k$ normal matrix is at least two, so
$k-n\ge2$.  Sharpness follows from Theorem~\ref{thm:parity-lift} below.
\end{proof}

\begin{theorem}[Parity lift of a central polygon]
\label{thm:parity-lift}
Let $n\ge2$ and let $k\ge n+2$ be even.  There is a bounded convex projective
$n$-polytope with exactly $k$ facets and a strict primitive open family of
$k$-periodic billiard trajectories.  Its canonical lifted mirror product is
$+I_{n+1}$, its witness makes exactly one full vector turn, and its mirror group
is infinite.
\end{theorem}

\begin{proof}
Put $k=2r$ and consider the $(k-2)$-dimensional space
\[
 S=\left\{s\in\mathbb R^k:
       \sum_{i\ \mathrm{even}}s_i=0,\quad
       \sum_{i\ \mathrm{odd}}s_i=0\right\}.               \tag{3}
\]
Let $v_i=(\cos(2\pi i/k),\sin(2\pi i/k))$ and define
\[
 F:\mathbb R^2\longrightarrow S,\qquad F(q)_i=\langle v_i,q\rangle .
\]
Indeed $F(q)$ lies in $S$: the even $v_i$ form a centred regular $r$-gon, and
the odd ones are its rotation by $2\pi/k$, so both parity sums vanish.  The map is
injective.  Since $2\le n\le k-2$, choose an $n$-plane
\[
                         F(\mathbb R^2)\subset W\subset S.  \tag{4}
\]

Let $N_i\in W^*$ be the restriction of the $i$-th coordinate functional and set
\[
                         \Omega=\{w\in W:N_i(w)<1\ \forall i\}.
\]
The coordinate restrictions span $W^*$.  Moreover, (3) gives
\[
             \sum_{i\ \mathrm{even}}N_i=0,
             \qquad \sum_{i\ \mathrm{odd}}N_i=0.          \tag{5}
\]
Thus $\sum_iN_i=0$ with all coefficients positive.  If $y$ were a recession
direction, every $N_i(y)\le0$; summing and using (5) would force every value to
vanish, and the spanning property would give $y=0$.  Hence $\Omega$ is bounded.

Every inequality is a genuine facet.  The point $F(v_i)$ satisfies
\[
 N_i(F(v_i))=1,
 \qquad N_j(F(v_i))=\langle v_j,v_i\rangle<1\quad(j\ne i),
\]
so it lies in the relative interior of facet $i$.  The first equation in (5)
minus the second is exactly the alternating relation (1), and
Proposition~\ref{prop:common-centre-product} gives product $+I_{n+1}$.

The plane $F(\mathbb R^2)$ cuts $\Omega$ in
\[
                         \langle v_i,q\rangle<1\quad(0\le i<k),
\]
the regular $k$-gon circumscribed about the unit circle.  Every common-centre
homology preserves the homogeneous span of this plane and restricts to the
centrally projective reflection in the corresponding side.  The standard
midpoint orbit therefore visits the $k$ ambient facets in cyclic order.  At each
impact all competing inequalities are strict, so it is the unique first hit and
lies in the relative interior of the ambient facet.  The central regular even
polygon witness has positive return multiplier and exactly one full vector turn.
The ambient openness criterion now gives a primitive open family, primitivity
following from the $k$ distinct facets.

Finally, by (2), for $i\ne j$
\[
 H_jH_i=\begin{pmatrix}I_n&0\\2(N_j-N_i)&1\end{pmatrix}
\]
is a nonidentity unipotent and has infinite order.  An independent sign and
dimension audit is provided by
\texttt{verificadores/}\allowbreak\texttt{parity\_lift.py}; the proof above,
rather than that numerical audit, establishes the result for all parameters.
\end{proof}

\begin{corollary}[Complete existence table]
\label{cor:complete-existence-table}
Under the convention that the period-$k$ witness visits the $k$ distinct facets
of its polytope, the determinant obstruction is the only obstruction to
existence:
\begin{enumerate}
\item if $n\ge2$ is even, every $k\ge n+1$ occurs;
\item if $n\ge3$ is odd, exactly the even $k\ge n+1$ occur.
\end{enumerate}
Every example can be chosen with a strict primitive witness and hence with an open
family of periodic trajectories.
\end{corollary}

\begin{proof}
Under the stated convention, a bounded $n$-polytope has at least $n+1$ facets,
so $k\ge n+1$ is necessary.  For every $n\ge2$, the negative-monodromy
construction gives all periods
\[
                         k=n+1+2q,\qquad q\ge0.             \tag{6}
\]
If $n$ is even, (6) is every odd $k\ge n+1$, while
Theorem~\ref{thm:parity-lift} gives every even $k\ge n+2$.  If $n$ is odd, (6)
is every even $k\ge n+1$, and the determinant parity theorem excludes odd $k$.
\end{proof}

\section{A metric obstruction for positive monodromy}
\begin{proposition}[Definite positive monodromy costs twice the rank]
\label{prop:finite-positive-twice-rank}
Let a bounded projective polytope in $\mathbb P(V)$, $\dim V=m$, have distinct
facet homologies whose canonical lifted product is $+I_V$.  If the homologies
preserve a common positive definite quadratic form, then their number satisfies
\[
                              k\ge2m.
\]
The same conclusion holds, in particular, if their mirror group is finite.
Consequently, every positive-monodromy example with $k<2(n+1)$ in projective
dimension $n$ has infinite mirror group and preserves no positive definite conic.
\end{proposition}

\begin{proof}
Use the invariant inner product and write the factors as orthogonal reflections
$t_i=s_{\alpha_i}$.  Put
\[
 w_i=t_i\cdots t_1,\qquad d_i=\operatorname{codim}\operatorname{Fix}(w_i),
 \qquad S_i=\operatorname{span}(\alpha_1,\ldots,\alpha_i).
\]
For an orthogonal map $w$, its moved space
$\operatorname{Mov}(w)=\operatorname{im}(w-I)$ is
$\operatorname{Fix}(w)^\perp$.  Multiplication by one reflection changes
$d=\dim\operatorname{Mov}(w)$ by exactly one.  More precisely,
\[
 d_i=d_{i-1}-1
 \quad\Longleftrightarrow\quad
 \alpha_i\in\operatorname{Mov}(w_{i-1}).               \tag{1}
\]
Indeed, if $\alpha_i\notin\operatorname{Mov}(w_{i-1})$ and
$t_iw_{i-1}x=x$, then
$(w_{i-1}-I)x$ is a multiple of $\alpha_i$ lying in the moved space, hence is
zero; thus
$\operatorname{Fix}(t_iw_{i-1})=\operatorname{Fix}(w_{i-1})\cap\alpha_i^\perp$
and $d_i=d_{i-1}+1$.  If $\alpha_i$ belongs to the moved space, then
$\operatorname{Fix}(w_{i-1})\subset\alpha_i^\perp$ is fixed by $t_iw_{i-1}$;
the rank-one perturbation and the determinant parity force
$d_i=d_{i-1}-1$.  This proves (1).

Every preceding reflection fixes $S_{i-1}^\perp$, so
$\operatorname{Mov}(w_{i-1})\subset S_{i-1}$.  Therefore a decreasing step in
(1) does not enlarge the span of the roots; only an increasing step can enlarge
it, and then by at most one.  Since $w_0=w_k=I$, the integer path
$d_0,d_1,\ldots,d_k$ starts and ends at zero with steps $\pm1$.  It has exactly
$k/2$ increasing steps.  Consequently
\[
                  \dim S_k\le k/2.                     \tag{2}
\]

The facet covectors span $V^*$: otherwise their common kernel would contain a
nonzero vector, and the homogeneous cone of the bounded polytope would contain a
line.  Under the invariant inner product this is equivalent to
$S_k=V$, so (2) gives $m\le k/2$.

Finally, a finite mirror group preserves a positive definite form obtained by
averaging any inner product over the group.  Each facet homology is then an
orthogonal reflection for that form, so the first part applies.
\end{proof}

\begin{remark}
Distinctness is required dynamically for primitive period but not in the group
argument.  Reducibility does not weaken the bound because the root-span argument
uses the full rank directly.  Degenerate and indefinite invariant forms are not
covered; the even central polygons used in the positive join preserve a degenerate
form, exactly as the proposition predicts for their short periods.
\end{remark}

\section{The metric period spectrum}

Call a configuration \emph{metric} if all facet homologies preserve one positive
definite quadratic form. Projectively, this is exactly a spherical billiard in a
bounded affine chart.

\begin{theorem}[The full negative-monodromy metric ray]\label{thm:metric-ray}
For every vector rank $m\ge3$ and every $q\ge0$, there is a bounded metric
projective billiard in $\mathbb{RP}^{m-1}$ with a primitive open family of period
\[
                              k=m+2q
\]
and canonical lifted product $-I_m$.
\end{theorem}

\begin{proof}
Start in rank three with the odd cyclic atom of period $p=3+2q$ used in
Lemma~\ref{lem:odd-cyclic-one-turn}. The displayed form
$Q=\operatorname{diag}(1,1,\lambda v)$ is positive definite, and that lemma proves
that the developed witness goes from a lift $x$ to $-x$ in exactly one turn.

Apply Proposition~\ref{prop:rank-one-suspension} $m-3$ times. If the old mirrors
preserve $Q>0$, their lifts $H_i\oplus1$ and the new base reflection
$J=I\oplus(-1)$ preserve $Q\oplus[1]>0$. Hence each suspension preserves
metricity as well as boundedness, strictness, primitivity, one turn and product
$-I$. The final rank and period are
\[
                         3+(m-3)=m,
       \qquad (3+2q)+(m-3)=m+2q .
\]
\end{proof}

Combining Theorem~\ref{thm:metric-ray} with
Proposition~\ref{prop:finite-positive-twice-rank} gives the current exact table:
\[
\begin{array}{c|l}
\text{range} & \text{metric status}\\ \hline
k=m & \text{exists and is the orthogonal simplex}\\
k=m+2q,\ q\ge1 & \text{exists with product }-I_m\\
\text{$m$ odd, $k$ even, }m<k<2m & \text{impossible}\\
\text{$m$ odd, $k$ even, }k\ge2m & \text{open}
\end{array}
\]
When $m$ is even, determinants force $k$ even, so the second row supplies every
allowed period. Metric existence is therefore completely determined in every odd
projective dimension. In even projective dimension the first unresolved pair is
$(m,k)=(3,6)$; more generally the unresolved branch is precisely odd $m$,
positive monodromy and even $k\ge2m$.
\section{Conclusion}
Under the stated all-facets convention, the parity lift and negative
monodromy constructions complete the non-metric period spectrum.  The metric
part is complete in odd projective dimension, while the only unresolved branch
is the one identified in the final table.  All finite-dimensional examples used
in the constructions are accompanied by exact rational certificates.
\newpage
\section*{Reproducibility and computational disclosure}
The general constructions are analytic. Supplementary exact-arithmetic scripts
check the rational atoms and representative joins, suspensions and parity lifts.
No floating-point result is used as a premise of a theorem. Research and manuscript development for all three papers used AI assistance from early
July through late August 2026 (approximate). The models used were OpenAI's GPT-6 Astra,
GPT-5.6 Sol, and GPT-5.6 Luna, and Anthropic's Fable 5 and Opus 5, accessed through
ChatGPT Work, ChatGPT Codex, and Claude Code. No API was used. The assistance supported
research, manuscript drafting, development and checking of verification scripts, and revision.
Jorge Lucas Gonz\'alez reviewed the manuscripts and assumes full responsibility; no AI
system is an author.
\par\medskip
\noindent\textbf{Competing interests:} The author declares none.
\par\medskip
\begin{thebibliography}{99}
\footnotesize
\enlargethispage{2\baselineskip}
\interlinepenalty=10000
\bibitem{tab} S.~Tabachnikov, \emph{Introducing projective billiards},
Ergodic Theory Dynam. Systems \textbf{17} (1997), 957--976,
\href{https://doi.org/10.1017/S0143385797086239}{doi:10.1017/S0143385797086239}.
\bibitem{fie-ex} C.~Fierobe, \emph{Examples of projective billiards with open
sets of periodic orbits}, Discrete Contin. Dyn. Syst. \textbf{44} (2024),
3287--3301, \href{https://doi.org/10.3934/dcds.2024059}{doi:10.3934/dcds.2024059}.
\bibitem{opb} M.~Bialy, C.~Fierobe, A.~Glutsyuk, M.~Levi, A.~Plakhov and
S.~Tabachnikov, \emph{Open Problems on Billiards and Geometric Optics},
Arnold Math. J. \textbf{8}(3--4) (2022), 411--422,
\href{https://doi.org/10.1007/s40598-022-00198-y}{doi:10.1007/s40598-022-00198-y}.
\bibitem{fie-caustics} C.~Fierobe, \emph{Only quadrics have pseudo-caustics ---
on caustics of Riemannian, pseudo-Euclidean and projective billiards in higher
dimensions}, Math. Z. \textbf{312} (2026), Art. 67,
\href{https://doi.org/10.1007/s00209-026-03939-3}{doi:10.1007/s00209-026-03939-3};
\href{https://arxiv.org/abs/2306.06757}{arXiv:2306.06757}.
\end{thebibliography}
\end{document}
